\documentclass[11pt,a4paper]{article}
\usepackage[margin=20mm,top=19mm,bottom=20mm]{geometry}
\usepackage{fontspec,amsmath,amssymb,xcolor,enumitem,fancyhdr}
\setmainfont{NotoSansKR-Regular.ttf}[Path=../../../common/assets/fonts/,BoldFont=NotoSansKR-Semibold.ttf]
\setmonofont{Latin Modern Mono}
\XeTeXlinebreaklocale "ko"
\XeTeXlinebreakskip=0pt plus 1pt
\definecolor{ink}{HTML}{203B50}
\definecolor{rulegray}{HTML}{D2D9DF}
\pagestyle{fancy}\fancyhf{}
\fancyhead[L]{\small 머신러닝 개념 확인}\fancyhead[R]{\small W3}
\fancyfoot[C]{\small\thepage\ / 4}
\renewcommand{\headrulewidth}{0.3pt}
\setlength{\headheight}{15pt}
\setlength{\parindent}{0pt}\setlength{\parskip}{5pt}
\linespread{1.2}
\newcommand{\parttitle}[2]{{\large\bfseries\color{ink} #1}\par{\small #2}\par\vspace{5mm}}
\newcommand{\question}[2]{\textbf{#1.} #2\par}
\newcommand{\answerlines}[1]{\foreachline{#1}}
\newcount\linecount
\newcommand{\foreachline}[1]{\linecount=0\loop\ifnum\linecount<#1\vspace{8mm}\par{\color{rulegray}\hrule}\advance\linecount by1\repeat}
\begin{document}

{\LARGE\bfseries\color{ink} W3 · 분류와 모델의 일반화}\par
{\small 로지스틱 회귀 · 결정 경계 · 과적합 · 정규화}\par
\vspace{3mm}이름\enspace\rule{43mm}{0.3pt}\hfill 날짜\enspace\rule{35mm}{0.3pt}\par\vspace{3mm}
\parttitle{1. 개념 확인}{각 문항에서 가장 적절한 답 하나를 고르세요.}
\question{1}{메일을 스팸과 정상으로 구분하는 작업은 무엇에 해당할까요?}
\begin{enumerate}[label=(\arabic*),leftmargin=9mm,itemsep=0pt,topsep=2pt,parsep=0pt]
\item 연속적인 집값 예측
\item 입력 단위 변환
\item 이진 분류
\item 데이터 개수 계산
\end{enumerate}\vfill
\question{2}{시그모이드 함수가 유한한 입력값을 받아 출력하는 값의 범위는 무엇일까요?}
\begin{enumerate}[label=(\arabic*),leftmargin=9mm,itemsep=0pt,topsep=2pt,parsep=0pt]
\item $-1$보다 작다
\item $0$과 $1$ 사이이다
\item 항상 $1$보다 크다
\item 어떤 입력에서도 정확히 $0$이다
\end{enumerate}\vfill
\question{3}{양성 클래스의 예측 확률이 0.8이고, 0.5 이상이면 양성으로 분류합니다. 결과는 무엇일까요?}
\begin{enumerate}[label=(\arabic*),leftmargin=9mm,itemsep=0pt,topsep=2pt,parsep=0pt]
\item 음성
\item 확률을 계산할 수 없다
\item 실제 정답이 반드시 양성이다
\item 모델은 양성으로 분류한다
\end{enumerate}\vfill
\question{4}{양성으로 판정하는 기준을 0.5에서 0.8로 높였습니다. 같은 예측 확률들을 사용할 때 무엇이 맞을까요?}
\begin{enumerate}[label=(\arabic*),leftmargin=9mm,itemsep=0pt,topsep=2pt,parsep=0pt]
\item 양성으로 판정되는 사례 수는 늘지 않는다
\item 모든 예측 확률이 자동으로 커진다
\item 정확도가 반드시 높아진다
\item 모든 실제 정답이 바뀐다
\end{enumerate}\vfill
\question{5}{결정 경계는 무엇을 뜻할까요?}
\begin{enumerate}[label=(\arabic*),leftmargin=9mm,itemsep=0pt,topsep=2pt,parsep=0pt]
\item 데이터를 저장한 폴더의 경계
\item 학습 시간이 끝나는 시점
\item 입력 공간에서 모델의 분류 결과가 바뀌는 경계
\item 정답을 무작위로 나누는 선
\end{enumerate}\vfill
\newpage
\parttitle{1. 개념 확인 · 계속}{각 문항에서 가장 적절한 답 하나를 고르세요.}
\question{6}{실제 정답이 1인 사례에 대해 양성 확률을 0.01로 예측했습니다. 로지스틱 손실은 어떤 예측에 비해 더 클까요?}
\begin{enumerate}[label=(\arabic*),leftmargin=9mm,itemsep=0pt,topsep=2pt,parsep=0pt]
\item 확률을 0.001로 예측한 경우
\item 확률을 0.9로 예측한 경우
\item 동일하게 0.01로 예측한 경우
\item 예측 확률과 관계없이 손실은 같다
\end{enumerate}\vfill
\question{7}{학습 데이터에서는 매우 잘 맞지만, 처음 보는 데이터에서는 자주 틀립니다. 무엇을 의심할 수 있을까요?}
\begin{enumerate}[label=(\arabic*),leftmargin=9mm,itemsep=0pt,topsep=2pt,parsep=0pt]
\item 반드시 완벽한 일반화
\item 데이터가 너무 적게 저장된 파일 형식
\item 분류 문제를 풀 수 없다는 증거
\item 과적합
\end{enumerate}\vfill
\question{8}{정규화가 과적합을 줄이는 데 도움을 주는 방식은 무엇일까요?}
\begin{enumerate}[label=(\arabic*),leftmargin=9mm,itemsep=0pt,topsep=2pt,parsep=0pt]
\item 지나치게 큰 가중치에 비용을 부과해 복잡한 맞춤을 억제한다
\item 학습 정답을 모두 삭제한다
\item 오차가 큰 예측에 보상을 준다
\item 모든 모델을 완벽하게 만든다
\end{enumerate}\vfill
\question{9}{다른 조건이 같을 때 정규화 강도 $\lambda$를 너무 크게 설정하면 어떤 일이 생길 수 있을까요?}
\begin{enumerate}[label=(\arabic*),leftmargin=9mm,itemsep=0pt,topsep=2pt,parsep=0pt]
\item 반드시 모든 오차가 사라진다
\item 입력 특성이 자동으로 늘어난다
\item 가중치가 과하게 억제되어 필요한 패턴도 못 배울 수 있다
\item 실제 정답의 개수가 늘어난다
\end{enumerate}\vfill
\question{10}{분류 모델이 학습을 잘했는지 판단하는 방법으로 가장 적절한 것은 무엇일까요?}
\begin{enumerate}[label=(\arabic*),leftmargin=9mm,itemsep=0pt,topsep=2pt,parsep=0pt]
\item 학습 데이터 정확도만 보고 결론을 낸다
\item 학습에 쓰지 않은 데이터에서도 성능을 확인한다
\item 가중치가 크면 무조건 좋다고 판단한다
\item 예측 결과를 모두 같은 클래스로 바꾼다
\end{enumerate}\vfill
\newpage
\parttitle{2. 머신러닝 사이언티스트의 점검}{11번 · 의사코드를 읽고, 모델이 의도대로 작동하지 않는 원인을 찾으세요.}
스팸 메일이면 y=1, 정상이면 y=0입니다. 모델이 매우 자신 있게 틀린 답을 내놓을수록 손실이 더 작아지는 현상이 나타납니다. 다른 계산과 학습률은 적절하며, 모든 p는 0보다 크고 1보다 작다고 가정합니다. log는 자연로그입니다.
\par\vspace{3mm}
{\small\ttfamily\begin{tabular}{r l}
1 & \hspace*{0.0mm}z = dot(x, w) + b \\
2 & \hspace*{0.0mm}p = sigmoid(z) \\
3 & \hspace*{0.0mm}loss = y * log(p) + (1 - y) * log(1 - p) \\
4 & \hspace*{0.0mm}update\_parameters\_to\_minimize(loss) \\
\end{tabular}}\par\vspace{5mm}
\textbf{11(1)} 손실을 정의한 줄에서 빠진 것이 무엇인지 찾고, 그 줄을 한 줄로 고쳐 쓰세요.
\answerlines{5}
\vspace{6mm}
\textbf{11(2)} 실제 정답이 y=1일 때 p=0.9와 p=0.1을 비교해, 원래 코드가 왜 틀린 예측을 선호하는지 설명하세요. 참고로 log(0.9)는 약 -0.105, log(0.1)은 약 -2.303입니다.
\answerlines{5}
\vspace{6mm}
\newpage
\parttitle{3. 열린 질문}{12번 · 정해진 정답은 없습니다. 제안한 방법과 그 이유가 드러나면 충분합니다.}
문장을 숫자 목록으로 바꾸는 도구와, 소량의 긍정·부정 표시가 붙은 리뷰가 있습니다. 이 자료로 처음 보는 리뷰의 감정을 분류하면서, 리뷰들 사이의 새로운 패턴도 발견하고 싶습니다. 어떤 순서로 실험하겠습니까?
\par\vspace{3mm}
이름표가 있는 리뷰와 없는 리뷰를 각각 어떻게 활용할지, 결과를 어떻게 확인할지, 모델이 확신하지 못하는 리뷰는 어떻게 다룰지 자유롭게 설명하세요. 배운 방법을 조합하거나 새로운 아이디어를 제안해도 됩니다.
\par\vspace{3mm}
{\small 수식이나 코드를 반드시 쓸 필요는 없습니다. 글, 그림, 간단한 순서도로 답해도 좋습니다.}
\answerlines{15}
\vfill{\footnotesize 이 문제지는 Machine Learning Specialization Course 1의 해당 주차 개념을 바탕으로 새로 작성한 학습용 자료입니다.}
\end{document}